% =========================
% Resume — Centered Header + Extra Section Spacing
% =========================
\documentclass[10.5pt]{article}

\usepackage[margin=0.75in]{geometry}
\usepackage[hidelinks]{hyperref}
\usepackage{enumitem}
\usepackage{tabularx}
\usepackage{fontawesome5}
\usepackage{titlesec}
\usepackage{xcolor}
\usepackage{setspace}

\pagenumbering{gobble}
\setlength{\parindent}{0pt}
\setlength{\parskip}{0pt}

% Tight bullets
\setlist[itemize]{leftmargin=*, nosep}

% Section formatting + MORE SPACE
\titleformat{\section}{\bfseries\large}{}{0pt}{}[\titlerule]
\titlespacing*{\section}{0pt}{14pt}{8pt} % ↑ increased spacing

% Helpers
\newcommand{\contactsep}{\hspace{10pt}|\hspace{10pt}}
\newcommand{\roleline}[2]{\textbf{#1}\hfill \textbf{#2}}
\newcommand{\orgline}[2]{\textit{#1}\hfill \textit{#2}}

\begin{document}

% ===== Centered Header =====
\begin{center}
    {\LARGE \textbf{STEVEN YANG}} \\[8pt]
    \faPhone\ (647)469-6927
    \contactsep
    \faEnvelope\ \href{mailto:steve.yang@mail.utoronto.ca}{steve.yang@mail.utoronto.ca}
    \contactsep
    \faGlobe\ \href{https://Sytion06.github.io}{Website}
    \contactsep
    \faLinkedin\ \href{https://www.linkedin.com/in/steven-yang-7760293a3/}{linkedin}
\end{center}

\section*{Summary}
Second-year Data Science Specialist and Computer Science major at the University of Toronto with strong foundations in software design, data analysis, and statistical modeling. Experienced in Java-based system design, data pipelines, and applied statistical analysis through academic and team projects. Seeking software engineering or data-focused internship/research opportunities.

% ===== Technical Skills =====
\section*{Technical Skills}
\begin{itemize}
  \item \textbf{Programming Languages:} Python, Java, C, JavaScript, HTML/CSS
  \item \textbf{Developer Tools:} Git/GitHub, Linux, VS Code, IntelliJ, Jupyter, LaTeX
  \item \textbf{Frameworks/Libraries:} NumPy, pandas, matplotlib, scikit-learn, SciPy, statesmodels
\end{itemize}

% ===== Education =====
\section*{Education}
\roleline{University of Toronto}{Sep. 2024 -- May 2028 (Expected)}
\begin{itemize}
  \item \textit{Honours Bachelor of Science}, Data Science Specialist, Computer Science Major
  \item cGPA: 4.0/4.0
  \item Relevant CourseWorks: Software Design, Data Structures and Analysis, Data Science, Probability and Statistics, \\ Linear Algebra, Multivariable Calculus
\end{itemize}

% ===== Projects =====
\section*{Projects}

\roleline{Question Bank Platform}{2026}
\orgline{Java, OpenAI API, FullStack}{}
\begin{itemize}
  \item Built a full-stack Java application to extract and structure math exam questions from PDFs into a searchable digital database.
  \item Integrated \textbf{OCR} and multimodal \textbf{OpenAI Vision models} to parse question text, diagrams, metadata, and page-level images..
  \item Designed an asynchronous processing pipeline with confidence scoring and manual review flags to improve extraction reliability.
  \item Implemented automated PDF segmentation, text cleaning, and structured JSON storage using Spring Boot + JPA.
\end{itemize}

\vspace{8pt}

\roleline{GroupFlow}{2025}
\orgline{Java, Swing, Clean Architecture, Git, JUnit, MongoDB}{}
\begin{itemize}
  \item Designed and implemented a collaborative task and group management application using \textbf{Clean Architecture}, separating entities, use cases, and UI layers.
  \item Built core features including group creation, role-based permissions, task assignment, and progress tracking.
  \item Improved code maintainability and testability through interface-driven design and unit testing with \textbf{JUnit}.
  \item Collaborated in a team environment using Git-based workflows, pull requests, and code reviews.
\end{itemize}

\vspace{8pt}

\roleline{Data Analysis Project}{2024}
\orgline{Python, pandas, NumPy, matplotlib, scikit-learn}{}
\begin{itemize}
  \item Analyzed real-world datasets in a group of 4 to explore relationships between variables using summary statistics and data visualizations.
  \item Performed data cleaning and preprocessing, handling missing values and inconsistent data formats.
  \item Applied statistical techniques such as logistic regression and train/test split to support conclusions.
  \item Presented findings through a 20-page slide deck.
\end{itemize}

% ===== Experience =====
% \section*{Experience}
% \roleline{ROLE TITLE}{DATES}
% \orgline{Organization}{City, Province}
% \begin{itemize}
%   \item Worked on \textbf{WHAT} using \textbf{TECH}, achieving \textbf{RESULT}.
%   \item Collaborated with \textbf{N} teammates using Git-based workflows.
% \end{itemize}

% ===== Activities / Awards =====
\section*{Awards}
\begin{itemize}
  \item \textbf{Dean’s List Scholar} \hfill 2025 \\
        University of Toronto

  \vspace{8pt}
  \item \textbf{Woodsworth College Scholarship Fund} \hfill 2026 \\
        Woodsworth College, University of Toronto
\end{itemize}

\end{document}
